% Proof notes for the main article; notation can be adapted by the integrator.
\section{Processed-component Potts tensors}

Fix a plane knot projection, one of its checkerboard graphs, an order of
its $n$ edges (the crossings), and a fixed integer $q\ge5$.  An active
vertex is incident with both a processed and an unprocessed edge.  A
processed component is a connected component of the spanning subgraph
formed by the processed edges after vertices of degree zero in that
subgraph have been discarded.  Its active vertices need not be consecutive
in any cyclic order.  Let $f_i$ be the number of active vertices after
edge $i$ and let $g$ be the maximum number of active vertices in one
processed component at any completed stage.

For a component $H$ with active set $B$, retain the function
\[
 F_H(\sigma_B)=\sum_{\sigma_{V(H)\setminus B}}
       \prod_{e=uv\in E(H)}W_e(\sigma_u,\sigma_v),
 \qquad W_e(a,b)=\begin{cases}-x^{-\epsilon_e},&a=b,\\1,&a\ne b.\end{cases}
\]
Here $x^2-(q-2)x+1=0$ and $\epsilon_e\in\{1,-1\}$ is the exponent of
the omitted-edge smoothing.  A completely closed component contributes
its scalar partition function.  The function of the entire processed
graph is exactly the product of these component functions.  This remains
true although unprocessed edges may later join those components: their
active spins are retained, and no interaction from an unprocessed edge
has yet been included.

Every $F_H$ is invariant under simultaneous permutation of the $q$ spin
names.  Represent an orbit by the restricted-growth word $P$ of its
active-spin equality pattern, with $k\le q$ blocks.  The stored value is
the aggregate
\[
 A_H(P)=(q)_k F_H(\sigma_P),\qquad
 (q)_k=q(q-1)\cdots(q-k+1).
\]
In particular, the coefficients of both integer coordinates of $A_H(P)$
are divisible by $(q)_k$.  The exact implementation checks the remainder
before every such division.  In the modular $q=5$ implementation these
orbit sizes are units modulo $2^{61}-1$.

\subsection{The merging identity}
Suppose an edge joins different processed components with equality
patterns $P,Q$, having $k,l$ blocks.  Their relative color data are a
partial bijection identifying $s$ blocks on one side with $s$ on the other.
The number of resulting colors is $m=k+l-s\le q$.  Nonidentified right
blocks receive successive fresh labels, so every joint equality orbit is
represented exactly once.  For each such relative alignment the aggregate
before its joining-edge weight is
\[
 \frac{A_H(P)}{(q)_k}\frac{A_J(Q)}{(q)_l}(q)_m.
\]
This follows by evaluating the two invariant functions at one labeled
representative and then multiplying by the size of its joint orbit.
Multiplication by the joining-edge weight and forgetting inactive endpoint
vertices complete the update.  Forgetting means restricting the equality
word, canonically renaming surviving blocks, and adding aggregates with
equal resulting words.  It does not involve division or a further factor.

The number of relative alignments of patterns with $k,l$ blocks is exactly
\[
 L_q(k,l)=\sum_{s=\max(0,k+l-q)}^{\min(k,l)}
           \binom{k}{s}\binom{l}{s}s!.
\]
It is at most $q!$: fix the left labels, and embed the right labels in
the $q$ names.  Canonical relative alignments are distinct representatives
among at most $(q)_l\le q!$ injections.  The bound is attained for
$k=l=q$.  Thus the $q=5$ merge has at most $120$ alignments per pair of
table entries, and the $q=6$ merge has at most $720$.
The orbit-cardinality identity
\[
 \sum_{\text{relative alignments}}(q)_m=(q)_k(q)_l
\]
is an independent counting check of the normalization.  The tests verify
it for every $k,l\le q\le6$, as well as the number of assignments in which
one prescribed left and one prescribed right block receive the same color.

\subsection{The component-width bound}
Put $S_q(t)=\sum_{k=0}^{\min(q,t)}\left\{{t\atop k}\right\}$.
A component with $t$ active vertices has at most $S_q(t)\le q^t$ table keys.
At a completed stage every component has at most $g$ active vertices, so
the sum of all represented keys is $O(nq^g)$.  Immediately before an edge
is processed, new endpoint vertices are introduced as singleton tensors.
The merged pre-forget bag of that edge has at most $g+2$ vertices: after
the edge it is one processed component with at most $g$ active vertices,
and only the at most two endpoint vertices can disappear at this step.
This argument also covers a loop and the closing of a whole component.

There are at most $q^{g+2}$ pairs of input entries at a merge, at most
$q!$ relative alignments per pair, and polynomial work in $g$ to form and
canonicalize a key.  Hence the transfer uses
\[
 n\,\operatorname{poly}(g+1)\,q!\,q^{g+2}
\]
coefficient-ring operations, in addition to polynomial preprocessing.
Logical key storage is $O(nq^g)$ at completed stages, or
$O(nq^{g+2})$ during a fused merge.  Two input tables and their replacement
may be simultaneously allocated.  The implementation's key budget counts
the sum of the represented replacement and unaffected tables, not Python
allocator overhead or a hard process-memory ceiling.

All exact coordinates have $O(n\log q)$ bits.  Indeed multiplication by
either $-x$ or $-x^{-1}$ increases the coefficient $\ell^1$ norm by at most
$q-1$, and summing at most $q^{|V|}$ spin assignments bounds that norm by
$q^{|V|}(q-1)^n$.  Orbit normalization only divides exact aggregates.
Products formed in a merge concern disjoint processed edge sets and have
the same global bound.  The normalization expression for the unknot also
has $O(n\log q)$ bits.  For fixed $q$ the bit bound is therefore
\[
 n^{O(1)}\operatorname{poly}(g+1)\,q^{g+2}.
\]
The same proof gives the weaker global-frontier bound with $g$ replaced
by $\max_i f_i$.  If $w_i$ is the number of projection edges cut by the
processed crossing set, every mixed black face has at least two cut-edge
incidences, and each cut projection edge has exactly one black side.
Consequently $2f_i\le w_i$, without any disk-boundary assumption.

The result is a quasi-polynomial bound for this exact Jones specialization
whenever $g=O(\log^2 n)$.  It remains a one-sided knot filter.  A value
equal to the unknot value is inconclusive, and no universal bound on $g$
for the implemented scan orders is asserted.

\subsection{Practical qualification}
The original greedy crossing order keeps its processed projection
connected.  Adjacent projection crossings share a black face, so its
processed Tait edges are also connected.  Consequently factorization has
little to separate on those orders and can add overhead.  It is most
useful when a supplied order creates independent processed components.
The benchmark intentionally includes every one of the first twelve
accepted diagrams and shuffled orders from one fixed seed.  Several
remain slower than the matching scan even after factorization; the
prototype is optional, not an unconditional replacement.

The $q=5$ specialization also has an infinite blind family among weaving
knots.  The reported $m=5,7,11$ cases of the closure of
$(\sigma_1\sigma_2^{-1})^m$ are inconclusive in both modular and exact
$q=5$ arithmetic.  Exact $q=6$ detects these cases.  Neither arithmetic
exactness nor processed-component factorization implies that a fixed
Jones specialization detects all nontrivial knots.
